\sscp{Bandicoots \it{(Peramelidae)}}{C-06-02}
Terms for bandicoots (\it{Peramelidae}) in \tsc{Seram-Tanimbar-Bomberai}
languages are only confidently known for four languages, given below.
\ili{Uruangnirin} \it{yawat kawer} ‘bandicoot’
and \ili{Onin} \it{yawar teroba} ‘bandicoot’ are literally ‘earth rat/mouse’,
in the same way the local \ili{Malay} term for bandicoot,
\it{tikus tanah}, is ‘earth rat/mouse’ (\it{yawat}
and \it{yawar} = ‘rat,
mouse’, \it{kawer} and \it{roba} = ‘soil, earth’).\footnote{
		The \ili{Uruangnirin}
		and \ili{Onin} terms may be calques on the local \ili{Malay} \it{tikus tanah}
		‘bandicoot’ = ‘earth rat/mouse’.
		However, the source of this local \ili{Malay} term may itself be
		a calque on an indigenous language of Papua.
		The semantics indicate that
		parallel innovation is also likely.}

Apart from terms explicitly glossed ‘bandicoot’ \citet[237]{SmitsVoorhoeve1992a}
have an entry headed ‘marsupial’.
Which \ili{Malay} terms were used to elicit the responses included here are unknown.
However, the \ili{Onin}, \ili{Sekar},
and \ili{Arguni} words are similar enough to \ili{Uruangnirin} \it{yawat
kawer} and \ili{Onin} \it{yawar teroba} to raise the suspicion
that they may also designate ‘bandicoot’.\footnote{
		\citet[237]{SmitsVoorhoeve1992a} also give \ili{Kowiai} \it{indor}
		and \ili{Irarutu} \it{yemoga}, \it{ʤemuge} under their entry for ‘marsupial’,
		both of which are independently verified to mean ‘cuscus’ (see
		\srf{C-06-01}).}
\\

\noindent{\wg\st{2pt}\begin{tabularx}{\linewidth}
{llll>{\raggedright\arraybackslash}X}
\lsptoprule
%Language	&	ISO	&	Term	&	Gloss	&	Source	\\	\midrule
%\endfirsthead
Language	&	ISO	&	Term	&	Gloss	&	Source	\\	\midrule
%\endhead
\ili{Kei} Kecil	&	kei	&	halee	&	mouselike 	&	\cite{Travis2019}	\\	\hhline{~}
					&			&				&	marsupial\footnotemark	&	\hp{\,}	\\
\ili{Kei} Besar	&	kei	&	‹helé›&	bandicoot	&	\cite[196]{Ellen1993}	\\
\ili{Uruangnirin}	&	urn	&	yawat kawer;	&	bandicoot;	&	\cite{Visser2023} (lit. `earth rat')	\\	\hhline{~}
	&		&	yawat	&	mouse, rat	&	\hp{\,}	\\
\ili{Onin}-1	&	oni	&	yawat	&	marsupial, mouse	&	SV \citeyear[237, 243]{SmitsVoorhoeve1992a}	\\
\ili{Onin}-2	&	oni	&	yawar teroba	&	bandicoot	&	\cite{Donohue2010} (lit. `earth rat')	\\
\ili{Sekar}	&	skz	&	kawat	&	marsupial	&	SV \citeyear[237]{SmitsVoorhoeve1992a}	\\
\ili{Arguni}	&	agf	&	rawake	&	marsupial	&	SV \citeyear[237]{SmitsVoorhoeve1992a} (C. met.?)	\\
\ili{Erokwanas}	&	erw	&	adobat	&	rat, marsupial	&	\cite{Donohue2010}	\\
\lspbottomrule
\end{tabularx}}
\footnotetext{The Indonesian gloss of \ili{Kei} \it{halee}
is \it{tikus berpundi} `pouched mouse/rat'.
It is given with the following
encyclopaedic information: ``Its snout is placed on a young
girl's breasts to make them fashionably small.''}

%\noindent{\wg\st{2pt}\begin{xltabular}{\linewidth}
%{lllll>{\raggedright\arraybackslash}X}
%\lsptoprule
%%Language	&	ISO	&	Term	&	Gloss	&	Etymon	&	Source	\\	\midrule
%%\endfirsthead
%Language	&	ISO	&	Term	&	Gloss	&	Etymon	&	Source	\\	\midrule
%\endhead
%
%\lspbottomrule
%\end{xltabular}}